\documentclass[10pt,a4paper]{amsart}
\usepackage{amsmath,amssymb,amsthm,graphicx}
% Retained minimal preamble for the Stacks Project chapter more-algebra.tex (More on Algebra).
% The chapter file pinned-original-source.tex carries only \input{preamble} on its first line, so no
% preamble text lives inside the chapter itself.  The Stacks Project uses one single shared preamble for
% every chapter; the snapshot stages it identically as sources/stacks-chow/preamble.tex and
% sources/stacks-etale-cohomology/preamble.tex.  The block below is copied mechanically, line for line,
% from sources/stacks-chow/preamble.tex lines 14-18 (the Xy-pic setup) and 154-217 (the native theorem
% declarations together with the mathematical macros).  Nothing was added, removed or redefined; the
% wrapper class and the amsmath/amssymb/amsthm/graphicx base packages are supplied by the assembler.
% Chapter: more-algebra.tex -- More on Algebra.
% For commutative diagrams we use Xy-pic
\usepackage[all]{xy}

% We use 2cell for 2-commutative diagrams.
\xyoption{2cell}

% Theorem environments.
%
\theoremstyle{plain}
\newtheorem{theorem}[subsection]{Theorem}
\newtheorem{proposition}[subsection]{Proposition}
\newtheorem{lemma}[subsection]{Lemma}

\theoremstyle{definition}
\newtheorem{definition}[subsection]{Definition}
\newtheorem{example}[subsection]{Example}
\newtheorem{exercise}[subsection]{Exercise}
\newtheorem{situation}[subsection]{Situation}

\theoremstyle{remark}
\newtheorem{remark}[subsection]{Remark}
\newtheorem{remarks}[subsection]{Remarks}

\numberwithin{equation}{subsection}

% Macros
%
\def\lim{\mathop{\mathrm{lim}}\nolimits}
\def\colim{\mathop{\mathrm{colim}}\nolimits}
\def\Spec{\mathop{\mathrm{Spec}}}
\def\Hom{\mathop{\mathrm{Hom}}\nolimits}
\def\Ext{\mathop{\mathrm{Ext}}\nolimits}
\def\SheafHom{\mathop{\mathcal{H}\!\mathit{om}}\nolimits}
\def\SheafExt{\mathop{\mathcal{E}\!\mathit{xt}}\nolimits}
\def\Sch{\mathit{Sch}}
\def\Mor{\mathop{\mathrm{Mor}}\nolimits}
\def\Ob{\mathop{\mathrm{Ob}}\nolimits}
\def\Sh{\mathop{\mathit{Sh}}\nolimits}
\def\NL{\mathop{N\!L}\nolimits}
\def\CH{\mathop{\mathrm{CH}}\nolimits}
\def\proetale{{pro\text{-}\acute{e}tale}}
\def\etale{{\acute{e}tale}}
\def\QCoh{\mathit{QCoh}}
\def\Ker{\mathop{\mathrm{Ker}}}
\def\Im{\mathop{\mathrm{Im}}}
\def\Coker{\mathop{\mathrm{Coker}}}
\def\Coim{\mathop{\mathrm{Coim}}}

% Boxtimes
%
\DeclareMathSymbol{\boxtimes}{\mathbin}{AMSa}{"02}

%
% Macros for moduli stacks/spaces
%
\def\QCohstack{\mathcal{QC}\!\mathit{oh}}
\def\Cohstack{\mathcal{C}\!\mathit{oh}}
\def\Spacesstack{\mathcal{S}\!\mathit{paces}}
\def\Quotfunctor{\mathrm{Quot}}
\def\Hilbfunctor{\mathrm{Hilb}}
\def\Curvesstack{\mathcal{C}\!\mathit{urves}}
\def\Polarizedstack{\mathcal{P}\!\mathit{olarized}}
\def\Complexesstack{\mathcal{C}\!\mathit{omplexes}}
% \Pic is the operator that assigns to X its picard group, usage \Pic(X)
% \Picardstack_{X/B} denotes the Picard stack of X over B
% \Picardfunctor_{X/B} denotes the Picard functor of X over B
\def\Pic{\mathop{\mathrm{Pic}}\nolimits}
\def\Picardstack{\mathcal{P}\!\mathit{ic}}
\def\Picardfunctor{\mathrm{Pic}}
\def\Deformationcategory{\mathcal{D}\!\mathit{ef}}

% --- appended 2026-10-12 (Chief): the project own definitions of the reference/slogan/history
% display environments, copied mechanically from the shared Stacks preamble lines 8-12
% (\usepackage{verbatim} plus the three \newenvironment lines, which the chapter preamble keeps
% verbatim from the project preamble). Nothing added, removed or redefined: these environments
% carry no mathematical content and the project itself typesets them as comments. Without these
% three lines a counted statement whose own text uses them fails with Environment slogan undefined.
\usepackage{verbatim}
\newenvironment{reference}{\comment}{\endcomment}
\newenvironment{slogan}{\comment}{\endcomment}
\newenvironment{history}{\comment}{\endcomment}
\usepackage{cleveref}
\crefname{originallemma}{original lemma}{original lemmas}
\Crefname{originallemma}{Original lemma}{Original lemmas}
\makeatletter
\newlabel{algebra-lemma-flat-going-down}{{00HS}{source}{}{sourcepointer-1}{}}
\newlabel{algebra-lemma-flat-going-down@cref}{{[originallemma][0][]00HS}{[1][0][]source}{}{}{}}
\makeatother
\crefname{originallemma}{original lemma}{original lemmas}
\Crefname{originallemma}{Original lemma}{Original lemmas}
\makeatletter
\newlabel{algebra-lemma-integral-going-up}{{00GU}{source}{}{sourcepointer-2}{}}
\newlabel{algebra-lemma-integral-going-up@cref}{{[originallemma][0][]00GU}{[1][0][]source}{}{}{}}
\makeatother
\crefname{originallemma}{original lemma}{original lemmas}
\Crefname{originallemma}{Original lemma}{Original lemmas}
\makeatletter
\newlabel{algebra-lemma-local-tensor-with-integral}{{092Y}{source}{}{sourcepointer-3}{}}
\newlabel{algebra-lemma-local-tensor-with-integral@cref}{{[originallemma][0][]092Y}{[1][0][]source}{}{}{}}
\makeatother
\crefname{originallemma}{original lemma}{original lemmas}
\Crefname{originallemma}{Original lemma}{Original lemmas}
\makeatletter
\newlabel{algebra-lemma-mop-up}{{04GJ}{source}{}{sourcepointer-4}{}}
\newlabel{algebra-lemma-mop-up@cref}{{[originallemma][0][]04GJ}{[1][0][]source}{}{}{}}
\makeatother
\crefname{originallemma}{original lemma}{original lemmas}
\Crefname{originallemma}{Original lemma}{Original lemmas}
\makeatletter
\newlabel{algebra-lemma-normal-goes-up}{{033C}{source}{}{sourcepointer-5}{}}
\newlabel{algebra-lemma-normal-goes-up@cref}{{[originallemma][0][]033C}{[1][0][]source}{}{}{}}
\makeatother
\crefname{originallemma}{original lemma}{original lemmas}
\Crefname{originallemma}{Original lemma}{Original lemmas}
\makeatletter
\newlabel{algebra-lemma-quotient-henselization}{{05WQ}{source}{}{sourcepointer-6}{}}
\newlabel{algebra-lemma-quotient-henselization@cref}{{[originallemma][0][]05WQ}{[1][0][]source}{}{}{}}
\makeatother
\crefname{originallemma}{original lemma}{original lemmas}
\Crefname{originallemma}{Original lemma}{Original lemmas}
\makeatletter
\newlabel{lemma-branches}{{31155}{source}{}{sourcepointer-7}{}}
\newlabel{lemma-branches@cref}{{[originallemma][0][]31155}{[1][0][]source}{}{}{}}
\makeatother
\crefname{originallemma}{original lemma}{original lemmas}
\Crefname{originallemma}{Original lemma}{Original lemmas}
\makeatletter
\newlabel{lemma-henselization-reduced}{{11067}{source}{}{sourcepointer-8}{}}
\newlabel{lemma-henselization-reduced@cref}{{[originallemma][0][]11067}{[1][0][]source}{}{}{}}
\makeatother
\title{A local ring is unibranch exactly when its henselization has a unique minimal prime}
\author{The Stacks Project Authors}
\date{stacks-project:more-algebra.tex:local-snapshot-20261012; GNU FDL 1.2 or later}
\begin{document}
\maketitle
\setcounter{subsection}{1}
\noindent\textit{Editorial note.} All delivered mathematical text is the author's own, copied from the pinned Stacks Project chapter file more-algebra.tex; no statement or proof step was rewritten, reordered or removed.\par\medskip
\noindent\textit{Editorial note.} The chapter file pinned-original-source.tex carries only \textbackslash{}input\{preamble\} on its first line, so no preamble text lives inside the chapter. The retained file sources/stacks-more-algebra/retained-preamble.tex is a mechanical, line-for-line copy of the project's single shared preamble as staged in the snapshot, together with the project's own definitions of the reference/slogan/history display environments, which the project itself typesets as comments; nothing was added, removed or redefined by this card and no package was added.\par\medskip
\noindent\textit{Editorial note.} The source URL is the source-level Stacks Project URL, as required when more than one card is produced from one source; the per-claim tag is carried by claim\_key.\par\medskip
\noindent\textit{Editorial note.} The counted statement span includes the author's own \textbackslash{}begin\{reference\} block citing the published source of the equivalence; it is delivered verbatim and produces no unresolved citation. The author writes that the result 'follows from' the branch lemma but gives a direct proof, which is what the counted proof span contains. The eight references inside the counted proof span are recorded as non-load-bearing navigation pointers.\par\medskip
\noindent\textit{Allowed standard prerequisites (editorial).}\begin{itemize}
\item Henselization of a local ring, its flatness, and the identity (A/p)\textasciicircum{}h = A\textasciicircum{}h/pA\textasciicircum{}h. used at proof lines 31269-31276. Is what identifies the unique minimal prime of A\textasciicircum{}h with the extension of that of A.
\item Unibranchness as locality of the integral closure of the reduced quotient. used at proof lines 31277-31278. Fixes the equivalent form of the conclusion that the direct proof establishes.
\item Integrality, going up, and flat base change of a chain of ring inclusions. used at proof lines 31279-31301. Produces the chain of inclusions over A\textasciicircum{}h whose maximal ideals give the contradiction.
\end{itemize}\medskip
\ifdefined\hypertarget\hypertarget{sourcepointer-1}{}\fi
\noindent\textit{Source reference.} Original lemma 00HS: focused-seeds/AGP/sources/stacks-snapshot-20261011/tags.txt, source lines 654--654. Cross-chapter navigation pointer into the Stacks Project; the standard fact it supplies is applied inside the counted proof span and the argument resting on it stays there.\par\medskip
\ifdefined\hypertarget\hypertarget{sourcepointer-2}{}\fi
\noindent\textit{Source reference.} Original lemma 00GU: focused-seeds/AGP/sources/stacks-snapshot-20261011/tags.txt, source lines 620--620. Cross-chapter navigation pointer into the Stacks Project; the standard fact it supplies is applied inside the counted proof span and the argument resting on it stays there.\par\medskip
\ifdefined\hypertarget\hypertarget{sourcepointer-3}{}\fi
\noindent\textit{Source reference.} Original lemma 092Y: focused-seeds/AGP/sources/stacks-snapshot-20261011/tags.txt, source lines 11496--11496. Cross-chapter navigation pointer into the Stacks Project; the standard fact it supplies is applied inside the counted proof span and the argument resting on it stays there.\par\medskip
\ifdefined\hypertarget\hypertarget{sourcepointer-4}{}\fi
\noindent\textit{Source reference.} Original lemma 04GJ: focused-seeds/AGP/sources/stacks-snapshot-20261011/tags.txt, source lines 5835--5835. Cross-chapter navigation pointer into the Stacks Project; the standard fact it supplies is applied inside the counted proof span and the argument resting on it stays there.\par\medskip
\ifdefined\hypertarget\hypertarget{sourcepointer-5}{}\fi
\noindent\textit{Source reference.} Original lemma 033C: focused-seeds/AGP/sources/stacks-snapshot-20261011/tags.txt, source lines 4051--4051. Cross-chapter navigation pointer into the Stacks Project; the standard fact it supplies is applied inside the counted proof span and the argument resting on it stays there.\par\medskip
\ifdefined\hypertarget\hypertarget{sourcepointer-6}{}\fi
\noindent\textit{Source reference.} Original lemma 05WQ: focused-seeds/AGP/sources/stacks-snapshot-20261011/tags.txt, source lines 7594--7594. Cross-chapter navigation pointer into the Stacks Project; the standard fact it supplies is applied inside the counted proof span and the argument resting on it stays there.\par\medskip
\ifdefined\hypertarget\hypertarget{sourcepointer-7}{}\fi
\noindent\textit{Source reference.} Original lemma 31155: sources/stacks-more-algebra/pinned-original-source.tex, source lines 31155--31155. Same-chapter navigation pointer; the statement it names is standard for the section and the argument that consumes it stays inside the counted proof span.\par\medskip
\ifdefined\hypertarget\hypertarget{sourcepointer-8}{}\fi
\noindent\textit{Source reference.} Original lemma 11067: sources/stacks-more-algebra/pinned-original-source.tex, source lines 11067--11067. Same-chapter navigation pointer; the statement it names is standard for the section and the argument that consumes it stays inside the counted proof span.\par\medskip
\begin{lemma}
\label{lemma-unibranch}
\begin{reference}
\cite[Chapter IV Proposition 18.6.12]{EGA4}
\end{reference}
Let $A$ be a local ring. Let $A^h$ be the henselization of $A$.
The following are equivalent
\begin{enumerate}
\item $A$ is unibranch, and
\item $A^h$ has a unique minimal prime.
\end{enumerate}
\end{lemma}

\begin{proof}
This follows from Lemma \ref{lemma-branches} but we will also give a
direct proof. Denote $\mathfrak m$ the maximal ideal of the ring $A$.
Recall that the residue field $\kappa = A/\mathfrak m$ is the same as
the residue field of $A^h$.

\medskip\noindent
Assume (2). Let $\mathfrak p^h$ be the unique minimal prime of
$A^h$. The flatness of $A \to A^h$ implies that
$\mathfrak p = A \cap \mathfrak p^h$ is the unique minimal
prime of $A$ (by going down, see
Algebra, Lemma \ref{algebra-lemma-flat-going-down}).
Also, since $A^h/\mathfrak pA^h = (A/\mathfrak p)^h$ (see
Algebra, Lemma \ref{algebra-lemma-quotient-henselization})
is reduced by Lemma \ref{lemma-henselization-reduced}
we see that $\mathfrak p^h = \mathfrak pA^h$.
Let $A'$ be the integral closure of $A/\mathfrak p$ in its fraction
field. We have to show that $A'$ is local.
Since $A \to A'$ is integral, every maximal ideal of $A'$ lies over
$\mathfrak m$ (by going up for integral ring maps, see
Algebra, Lemma \ref{algebra-lemma-integral-going-up}).
If $A'$ is not local, then we can find distinct maximal ideals
$\mathfrak m_1$, $\mathfrak m_2$. Choose elements $f_1, f_2 \in A'$
with $f_i \in \mathfrak m_i$ and $f_i \not \in \mathfrak m_{3 - i}$.
We find a finite subalgebra $B = A/\mathfrak p[f_1, f_2] \subset A'$
with distinct maximal ideals $B \cap \mathfrak m_i$, $i = 1, 2$.
Note that the inclusions
$$
A/\mathfrak p \subset B \subset \kappa(\mathfrak p)
$$
give, on tensoring with the flat ring map $A \to A^h$ the inclusions
$$
A^h/\mathfrak p^h \subset
B \otimes_A A^h \subset
\kappa(\mathfrak p) \otimes_A A^h \subset
\kappa(\mathfrak p^h)
$$
the last inclusion because
$\kappa(\mathfrak p) \otimes_A A^h =
\kappa(\mathfrak p) \otimes_{A/\mathfrak p} A^h/\mathfrak p^h$
is a localization of the domain $A^h/\mathfrak p^h$.
Note that $B \otimes_A \kappa$ has at least two maximal ideals
because $B/\mathfrak mB$ has two maximal ideals. Hence, as
$A^h$ is henselian we see that
$B \otimes_A A^h$ is a product of $\geq 2$ local rings, see
Algebra, Lemma \ref{algebra-lemma-mop-up}.
But we've just seen that $B \otimes_A A^h$ is a subring of a domain
and we get a contradiction.

\medskip\noindent
Assume (1). Let $\mathfrak p \subset A$ be the unique minimal
prime and let $A'$ be the integral closure of $A/\mathfrak p$
in its fraction field. Let $A \to B$ be a local map of local rings
inducing an isomorphism of residue fields which is a
localization of an \'etale $A$-algebra. In particular $\mathfrak m_B$
is the unique prime containing $\mathfrak m B$. Then $B' = A' \otimes_A B$
is integral over $B$ and the assumption that $A \to A'$ is local
implies that $B'$ is local
(Algebra, Lemma \ref{algebra-lemma-local-tensor-with-integral}).
On the other hand, $A' \to B'$ is the localization
of an \'etale ring map, hence $B'$ is normal, see
Algebra, Lemma \ref{algebra-lemma-normal-goes-up}.
Thus $B'$ is a (local) normal domain. Finally, we have
$$
B/\mathfrak pB \subset B \otimes_A \kappa(\mathfrak p)
= B' \otimes_{A'} (\text{fraction field of }A') \subset
\text{fraction field of }B'
$$
Hence $B/\mathfrak pB$ is a domain, which implies that $B$ has a unique
minimal prime (since by flatness of $A \to B$ these all have to lie
over $\mathfrak p$). Since $A^h$ is a filtered colimit of
the local rings $B$ it follows that $A^h$ has a unique minimal prime.
Namely, if $fg = 0$ in $A^h$ for some non-nilpotent elements
$f, g$, then we can find a $B$ as above containing both $f$ and $g$
which leads to a contradiction.
\end{proof}

\end{document}
